\section{Empirical implications}\label{sec:12_empirical}

The empirical results are predictions of the model, not estimates from takeover data. They use the benchmark
auction, the prior \(1/2\) and a slot in the band of Proposition~A.1, and they concern execution prices, preparation
decisions and winners. \paper{Sec.~8; Appendix~A.8}

\subsection{The domain and the observable order law}\label{sec:12_empirical-domain}

\paragraph{Full two-sided trade on \((r_S,h)\).} Let \(r_S\) be the strength at which a full sale earns the trading
cost, \(\rho\DT(r_S)m=k\); \Cref{sec:08_coexistence-thresholds} solves it with the root map. For \(r_S<r<h\), the
bounds of \Cref{sec:05_investor-bounds} give, under every candidate,
\[
K_H\ge\rho\DT(r)(1-M)=\rho\DT(r)m>k,\qquad K_L\ge\rho\DT(r)m>k,
\]
because \(1-M=m\) at the prior \(1/2\) and \(\DT\) rises in \(r\). So full orders are unique under both investor rules.
The interval contains \(r_1\), \(r_C\) and \(r_2\). Use the exact root for strict claims: the display
\(r_S\approx\reg{rS_B}\) lies below \(r_S\), and at \(r=\reg{rS_counter_r_B}\), above the display, the full sale
earns \(\rho\DT m-k=\reg{rS_counter_margin_B}<0\).

\paragraph{The order law.} Under full orders class \(H\) buys with probability \(b+\lambda\) and class \(L\) with
probability \(b\). At the prior \(1/2\), write \(\bar f=b+\lambda/2\). Then \(\Pr(X=\pm1)=\bar f\),
\(\Pr(X=0)=1-2\bar f\) and \(\E[X]=0\). The volume \(v=\Pr(X\ne0)=2b+\lambda\) counts nonzero orders per order
opportunity. The informed share of nonzero orders is
\[
\vartheta=\frac{\lambda}{2b+\lambda}=\frac{u}{2+u}=M-m=2M-1,
\]
where the last two forms use \(M=(1+u)/(2+u)\) and \(m=1/(2+u)\) at the prior \(1/2\). The share differs from the
presence probability: at the benchmark, \(v=\reg{vol_full_B}\), \(\vartheta=\reg{theta_B}\) and
\(\lambda=\reg{lambda_B}\).

\paragraph{Entry.} Write \(A_r(\mu)=t_L-t_0+\DT(r)\mu\); this is the average entry gain \(w(\mu)\) of
\Cref{sec:04_slot_beliefs-price} at strength \(r\). Let \(\chi_r=\1\{B_r(M)\ge c_H\}\). Only a purchase can recruit, and
\(B_r(M)\) falls strictly in \(r\) (\Cref{sec:06_main_result-ceiling}), so under rule C1 \(\chi_r=1\) exactly when
\(r\le r_C\). Entry after \((-1,0,1)\) is \((\rho,\rho,e_+)\) with \(e_+=\rho+(1-\rho)\chi_r\), so
\(\Ent=\rho+(1-\rho)\chi_r\bar f\) (\Cref{sec:06_main_result-entry}). From \(r_0\) to \(r_1\) entry rises by
\((1-\rho)\bar f=\reg{entry_jump_B}\). With the challenger's information held fixed, a stronger incumbent cannot raise
entry (\Cref{sec:07_controls-fixed}); the rise needs the change in information.

\subsection{The execution spread and its entry component}\label{sec:12_empirical-spread}
\paper{Prop.~A.11}

\begin{intuition}[A purchase moves two things]
A sale and a purchase move the class belief by the same amount in opposite directions. A purchase can also recruit a
high-cost challenger. That second effect enters the purchase price alone.
\end{intuition}

By \Cref{sec:04_slot_beliefs-price}, \(P(-1)=t_0+\rho A_r(m)\), \(P(0)=t_0+\rho A_r(1/2)\) and
\(P(1)=t_0+e_+A_r(M)\). Define the spread \(S=P(1)-P(-1)\), the purchase impact \(I_+=P(1)-P(0)\) and the sale
impact \(I_-=P(0)-P(-1)\).

\begin{resultbox}[The spread and its entry component \paper{Prop.~A.11}]
Let \(r_S<r<h\).
\begin{enumerate}[label=(\roman*),leftmargin=2em]
\item \(S=\rho\DT\vartheta+(1-\rho)\chi_rA_r(M)\): a class component and an entry component.
\item \(I_-=\rho\DT\vartheta/2\) and \(I_+=I_-+(1-\rho)\chi_rA_r(M)\). So \(I_+>I_-\) exactly when a purchase recruits.
\item With a hidden price, full orders remain unique, and \(S=\rho\DT\vartheta\).
\item On \((r_S,r_C]\), both components, both impacts and \(S\) rise strictly with \(r\). On \((r_C,h)\), the entry
component is zero, and the class component, both impacts and \(S\) rise strictly.
\item At \(r_C\), \(S\), \(P(1)\) and \(I_+\) fall by \((1-\rho)A_{r_C}(M)>0\); \(P(-1)\), \(P(0)\) and \(I_-\) are continuous.
\end{enumerate}
\end{resultbox}

\emph{Derivation.} Since \(e_+=\rho+(1-\rho)\chi_r\), \(S=\rho[A_r(M)-A_r(m)]+(1-\rho)\chi_rA_r(M)\). The function \(A_r\)
has slope \(\DT\), so \(A_r(M)-A_r(m)=\DT\vartheta\): this is (i). Since \(m+M=1\), \(\tfrac12-m=M-\tfrac12=\vartheta/2\),
which gives both impacts; and \(A_r(M)>0\), which gives the test in (ii). With a hidden price entry is \(\rho\) after
every order (\Cref{sec:07_controls-hidden}), so \(\chi_r\) drops out: (iii). At the benchmark \(t_L-t_0=\tfrac12-\tfrac1{8r}\)
and \(\DT\) rise in \(r\), so \(A_r(M)\) rises, while \(\rho\) and \(\vartheta\) do not move: (iv). Rule C1 gives
\(\chi_{r_C}=1\), and \(\chi_r=0\) above \(r_C\); every other term is continuous: (v). The test in (ii) needs no auction
values.

\paragraph{Benchmark check.} At \(r_1\), \(S=\reg{ba_class_r1_B}+\reg{ba_entry_r1_B}=\reg{ba_r1_B}\), so the entry
component is \(\reg{ba_share_r1_B}\) of the spread. The impacts are \(I_-=\reg{impact_sell_r1_B}\) and
\(I_+=\reg{impact_buy_r1_B}\). With a hidden price the spread would be \(\reg{ba_class_r1_B}\). At \(r_2\) it is
\(\reg{ba_r2_B}\), the class component alone, and at \(r_C\) it falls by \(\reg{ba_jump_rC_B}\). Under rule CM the value
at \(r_C\) can lie anywhere between the two one-sided limits.

\begin{caution}[This is not a quoted bid--ask spread]
The model has no quote stage. \(S\) is the difference between the execution prices of a unit purchase and a unit
sale, each conditional on its order. The sale impact uses the price after a zero order as its reference.
\end{caution}

\subsection{Identify the informed share from prices and counts}\label{sec:12_empirical-identification}
\paper{Prop.~A.12}

\paragraph{Counts alone.} Under full orders the order law \((v/2,1-v,v/2)\) and entry \(\rho+(1-\rho)\chi_rv/2\)
depend on \((\lambda,\beta)\) only through \(v\). Fix a ratio \(u'>0\) and solve \(v=2b'+\lambda'\), \(\lambda'=u'b'\):
\[
b'=\frac{v}{2+u'},\qquad \lambda'=\frac{vu'}{2+u'},\qquad \beta'=\frac{b'}{1-\lambda'}=\frac{v}{2+u'(1-v)}.
\]
This pair is unique, and it is a slot of the model exactly when \(\beta'\le1/3\), that is \(3v\le2+u'(1-v)\). In the
share \(\vartheta'=u'/(2+u')\), so that \(u'=2\vartheta'/(1-\vartheta')\), the right side is
\((2-2v\vartheta')/(1-\vartheta')\), and
\[
3v(1-\vartheta')\le2-2v\vartheta'\iff v\vartheta'\ge3v-2\iff\vartheta'\ge3-\frac2v.
\]
So every share from \(3-2/v\) up gives one slot with the same counts. The counts do not choose among them.

\paragraph{Incentives.} A slot with share \(\vartheta'\) has \(m'=(1-\vartheta')/2\). A sale never recruits, so under
full orders the low-class sale earns \(\rho\DT(r)m'-k\), which is nonnegative only if
\(\vartheta'\le\bar\vartheta(r)=1-2k/(\rho\DT(r))\). At \(r_1\), \(\bar\vartheta=1-\frac{1/25}{1/6}=\reg{theta_max_full_eq_r1_B}\).
So under rule IM the counts and the incentives at \(r_1\) allow
\(\vartheta'\in[3-2/v,\ \bar\vartheta(r_1)]\approx[\reg{theta_min_counts_B},\ \reg{theta_max_full_eq_r1_B}]\). At the
upper end the full sale earns exactly zero. A lower sale rate would raise \(\mu(-1)\) and make the sale strictly
profitable, so no mixture replaces the full sale. Rule I0 excludes the full sale there, and that slot then has no
equilibrium at \(r_1\).

\paragraph{The band.} Strict recruitment at \(r_1\) and strict exclusion at \(r_2\) are \(\tau(r_1)<M'<\tau(r_2)\), that
is, with \(\vartheta'=2M'-1\),
\[
2\tau(r_1)-1=\frac{41}{100}<\vartheta'<\frac{1343}{2902}=2\tau(r_2)-1.
\]
Every share in this band exceeds \(3-2/v\), so it has a slot with the benchmark counts and the benchmark entry at
\(r_1\) and \(r_2\). Full orders are unique at both strengths, because the weakest correct-order bound at \(r_1\),
\(\rho\DT(r_1)[1-\tau(r_2)]-k=\reg{margin_obs_band_incentive_B}\), is positive, and \(\DT(r_2)>\DT(r_1)\). For example, \(u'=8/5\) gives
\(\vartheta'=4/9=\reg{alt_theta_B}\) and the same entry, but a spread of \(\reg{alt_ba_r1_B}\) at \(r_1\), against
\(\reg{ba_r1_B}\). Prices separate slots that counts cannot.

\begin{resultbox}[Recover the slot from prices \paper{Prop.~A.12(d)}]
Let \(r_S<r<h\). Given \(\rho=\Pr(\text{preparation}\mid X\ne1)\) and the spread \(\DT(r)\), the sale impact and the
volume identify the slot:
\[
\vartheta=\frac{2I_-}{\rho\DT(r)},\qquad \lambda=v\vartheta,\qquad \beta=\frac{v(1-\vartheta)}{2(1-v\vartheta)}.
\]
\end{resultbox}
The first formula inverts \(I_-=\rho\DT\vartheta/2\); the second uses \(\vartheta=\lambda/v\); the third uses
\(b=v(1-\vartheta)/2\) and \(\beta=b/(1-\lambda)\). The floor is observable because entry after a sale and after zero is
\(\rho\). The target-payoff spread \(\DT(r)\) remains a maintained input of the auction model.

\paragraph{Counts under purchases only.} On \([r_F,r_N]\), form F of \Cref{sec:08_coexistence-set} has class \(H\)
buying with probability one and class \(L\) abstaining. Then \(\Pr(X=1)=b+\lambda/2\), \(\Pr(X=-1)=b\),
\(\E[X]=\lambda/2\) and \(v=2b+\lambda/2\), so \(\lambda=2\E[X]\) and \(b=(v-\E[X])/2\): the counts identify the slot.
On branch R class \(H\) buys with probability \(\sigma_H(1)<1\), and the same algebra identifies only the purchase rate
\(s_+=\lambda\sigma_H(1)=2\E[X]\) and \(b\), not \(\lambda\).

\subsection{Orders predict the winner after entry becomes flat}\label{sec:12_empirical-winner}
\paper{Prop.~A.13}

Let \(A\) be the event that the challenger prepares and \(W\) the event that it acquires the target. The cost is
independent of the class and of \(R\), so preparation leaves the class posterior unchanged:
\(\Pr(H\mid X=x,A)=\mu(x)\). A prepared high-class challenger always wins, because \(h>r\); a prepared low-class
challenger wins when \(R<\ell\), with probability \(\ell/r\). Thus
\[
w_x:=\Pr(W\mid X=x,A)=\mu(x)+[1-\mu(x)]\,\ell/r,\qquad \Pr(W\mid X=x)=e(x)\,w_x,
\]
and under full orders \(w_1-w_{-1}=(M-m)(1-\ell/r)=\vartheta(1-\ell/r)\).
The gap rises with \(r\), with derivative \(\vartheta\ell/r^2>0\). It is \(\reg{wgap_r1_B}\) at \(r_1\) and
\(\reg{wgap_r2_B}\) at \(r_2\).

Since \(\E[X]=0\), an indicator \(Y\) with conditional means \(y_x\) has \(\operatorname{Cov}(X,Y)=\E[XY]=\bar f(y_1-y_{-1})\).
For preparation \(y_1-y_{-1}=e_+-\rho\), so \(\operatorname{Cov}(X,\1_A)=(1-\rho)\chi_r\bar f\): \(\reg{cov_entry_r1_B}\)
at \(r_1\) and zero above \(r_C\). Above \(r_C\), \(e\equiv\rho\), so \(\operatorname{Cov}(X,\1_W)=\bar f\rho\vartheta(1-\ell/r)>0\),
which is \(\reg{cov_win_r2_B}\) at \(r_2\). Zero covariance between orders and entry does not mean that the price carries
no class information.

\subsection{Announcement returns and the zero order}\label{sec:12_empirical-announcements}
\paper{Prop.~A.14}

Add a public event after the price that reveals only whether the challenger prepared. With preparation the price
becomes \(t_0+A_r(\mu(x))\); without it, \(t_0\). Subtract \(P(x)=t_0+e(x)A_r(\mu(x))\):
\[
\Delta_x^1=[1-e(x)]A_r(\mu(x)),\qquad \Delta_x^0=-e(x)A_r(\mu(x)).
\]
The report of no preparation has probability \(1-e(x)\). When \(e(x)<1\), the mean change over both reports is
\(e(x)[1-e(x)]A_r-[1-e(x)]e(x)A_r=0\). When \(e(x)=1\), a report of preparation changes nothing.

A sample of preparation reports alone weights order \(x\) by \(\Pr(X=x\mid A)=\Pr(X=x)e(x)/\Ent\), so its mean change
is \(\sum_x\Pr(X=x\mid A)[1-e(x)]A_r(\mu(x))\). Every term is nonnegative. The posteriors average to the prior, so
some order has \(\mu(x)\le1/2\). At the benchmark \(B_r(1/2)<c_H\) for \(1<r<h\), so \(\tau(r)>1/2\ge\mu(x)\), and entry
at that order is \(\rho<1\). The order has positive probability, so the mean is positive. At \(r_1\) the changes after
a sale, zero and a purchase are \((\reg{evA_r1_sale_B},\ \reg{evA_r1_neutral_B},\ 0)\), with mean
\(\reg{evprep_mean_r1_B}\) in the sample of preparation reports. These are price changes; a percentage return divides
by the price before the event.

\paragraph{No order is bad news about entry.} The expected price before the order opportunity is
\(\bar P=\sum_x\Pr(X=x)P(x)\). Since \(A_r\) is affine and \(m+M=1\), \(\bar fA_r(m)+\bar fA_r(M)=2\bar fA_r(1/2)\), so
\begin{align*}
\bar P&=t_0+\rho\bigl[\bar fA_r(m)+(1-2\bar f)A_r(1/2)+\bar fA_r(M)\bigr]+(1-\rho)\chi_r\bar fA_r(M)\\
&=t_0+\rho A_r(1/2)+(1-\rho)\chi_r\bar fA_r(M)=P(0)+(1-\rho)\chi_r\bar fA_r(M).
\end{align*}
At \(r_1\), \(\bar P=\reg{pbar_r1_B}\), \(P(0)=\reg{p_neutral_r1_B}\), and \(P(0)-\bar P=\reg{pdrop_neutral_r1_B}\). The zero
order leaves the class posterior at \(1/2\) but removes the chance of recruitment after a purchase. This difference is
minus the access gain of \Cref{sec:11_welfare-access}: \(\bar P\) equals expected proceeds with access, and \(P(0)\)
equals expected proceeds with a hidden price. Above \(r_C\) the difference is zero.

\subsection{Two clusters at one intermediate strength}\label{sec:12_empirical-clusters}
\paper{Sec.~8; Prop.~3; Appendix~A.8}

For \(r_F<r<r_N\), Proposition~3 gives exactly two equilibria under rules IM and C1 (\Cref{sec:08_coexistence-set}).
In N both classes abstain, the three prices coincide and entry is \(\rho\). In F class \(H\) buys, class \(L\) abstains,
the mean order is \(\lambda/2=\reg{imb_buy_B}\), and every challenger prepares after a purchase.

F has a price pattern that a test can use. After a sale the posterior is \(1/2\), because \(f_H(-1)=f_L(-1)=b\). After zero,
\(f_H(0)=b_0\) and \(f_L(0)=b_0+\lambda\), with \(b_0=1-2b-\lambda\). With \(v_U=2b+\lambda/2\),
\[
\mu_0=\frac{b_0}{2b_0+\lambda}=\frac{1-2b-\lambda}{2-4b-\lambda},\qquad
\frac12-\mu_0=\frac{\lambda}{2(2-4b-\lambda)}=\frac{\lambda}{4(1-v_U)}.
\]
Both orders have entry \(\rho\), so \(P(-1)-P(0)=\rho\DT\lambda/[4(1-v_U)]>0\): a sale executes above the price after
zero. At \(r=\reg{str_r160}\) the difference is \(\reg{U_gap_r160_B}\). Rule CM adds the threshold form T as a third
tuple, so the two-cluster statement needs rule C1. The model does not say which equilibrium a process selects.

\subsection{Measurement: date the decision and retain zero orders}\label{sec:12_empirical-measurement}
\paper{Sec.~8}

\begin{enumerate}[leftmargin=2em]
\item Preparation is the start of substantive diligence, or a proposal that needs costly preparation; a final public
offer is a narrower, selected event.
\item Incumbent strength and the price must both be public before the further buyer's decision; a later price can
reflect anticipation of the bid.
\item The order is the first prespecified public order opportunity after the sale becomes public, such as a scheduled
call auction. An opportunity with no order, or an imbalance below a size fixed in advance, is a zero, and it stays in
the denominator of \(v\).
\item The announcement formulas need events that reveal preparation alone.
\end{enumerate}

\begin{caution}[The omitted zero changes the sample]
A sample of first executed trades discards the inactive opportunities. It cannot identify the volume, the price
effect of a zero order or the informed share. The share formula also needs \(\rho\) and \(\DT(r)\), so it is
conditional on the maintained auction model.
\end{caution}
